\documentclass[12pt]{article}
\usepackage{amsmath}
\usepackage{amssymb}
\usepackage{physics}
\usepackage{graphicx}
\usepackage{hyperref}

\newcommand{\ds}{\displaystyle}

\title{The Kinematics of Rigid Body Motion}
\author{}
\date{}

\begin{document}
\maketitle
 

% ============================================================
\section{Introduction}
% ============================================================

In this section, we will explore the kinematics of rigid body motion, focusing on the study
of the motion of objects that maintain a constant shape. Several comprehensive treatises have
been published on this topic; however, our goal is to delve deeper, and we will take the time
to derive each equation meticulously, even at the risk of appearing pedantic, to ensure
clarity in every notation.

We will cover the following topics:

\begin{itemize}
  \item \textbf{Basic Properties of Orthogonal Transformations:} We will outline some
        fundamental properties of orthogonal transformations within a vector space $V$ of
        dimension $n$, which is endowed with a positively defined scalar product. This will
        include a discussion on how these transformations preserve lengths and angles between
        vectors, a key aspect of their role in rigid body motion.
  \item \textbf{Rotations in Three Dimensions:} We will then delve into the significant case
        of rotations in three-dimensional space. Here, we will explore how rotational matrices
        are constructed and how they relate to the concept of angular velocity and rotation
        axes.
  \item \textbf{Infinitesimal Rotations:} From our understanding of three-dimensional
        rotations, we will introduce the concept of infinitesimal rotation. This involves
        small-angle approximations and the use of skew-symmetric matrices to represent these
        small rotations. We will discuss their properties and how they are used to derive the
        equations of rigid motion.
  \item \textbf{Transformation of Vectors Between Reference Frames:} Building upon
        infinitesimal rotations, we will derive the general formula for transforming vectors
        between fixed and rotating reference frames. This section will cover how vectors
        change when observed from different frames and the mathematical framework needed for
        these transformations.
  \item \textbf{Coriolis and Centrifugal Forces:} Finally, we will end with the derivation of
        the Coriolis and centrifugal forces. These fictitious forces arise due to the rotation
        of the reference frame and are crucial for understanding the dynamics in rotating
        systems, such as the Earth atmosphere and ocean currents.
\end{itemize}

% ============================================================
\section{Orthogonal Transformations}
% ============================================================

In this section we will review some basic properties of orthogonal transformations in a
vector space $V$ on the complex numbers $\mathbb{C}$ of dimension $n$, which is endowed with
a scalar product generalized for complex vectors:
\begin{equation}
  \ip{\vb{x}}{\vb{y}} = x_i\,\overline{y}_i
\end{equation}
where $\vb{x} = (x_1, x_2, \ldots, x_n)$ and $\vb{y} = (y_1, y_2, \ldots, y_n)$ are complex
vectors in $V$, and $\overline{y}$ denotes the complex conjugate of $y$. The scalar product
is positively defined, which means that it satisfies the following properties:
\begin{equation}
  \begin{gathered}
    \ip{\vb{x}}{\vb{y}} = \overline{\ip{\vb{y}}{\vb{x}}} \\
    \ip{\vb{x}}{\vb{x}} \geq 0 \\
    \ip{\vb{x}}{\vb{x}} = 0 \iff \vb{x} = \vb{0}
  \end{gathered}
\end{equation}

We will review some basic properties of orthogonal transformations, and we will consider only
orthogonal transformations with \textbf{real numbers} coefficients, as this is the physically
meaningful case for rigid body motion, although we will keep the scalar product generalized
for complex vectors.

\begin{itemize}
  \item \textbf{Definition}: Orthogonal transformations are linear transformations that
        preserve the length of vectors, and play a crucial role in the theory of rigid body
        motion. A linear transformation $\vb{A}$ of $V$ onto itself,
        \begin{equation}
          \vb{A}: V \to V
        \end{equation}
        is orthogonal if it preserves the length of a vector in both coordinate systems
        $\vb{x}$ and $\vb{x}'$ in $V$,
        \begin{equation}
          \ip{\vb{x}'}{\vb{x}'} = \ip{\vb{A}\vb{x}}{\vb{A}\vb{x}} = \ip{\vb{x}}{\vb{x}}
        \end{equation}
        for all $\vb{x}, \vb{x}' \in V$.

  \item \textbf{Symmetry and Inverse:} From the properties of matrices and scalar products,
        we can write the above equation as
        \begin{equation}
          \ip{\vb{A}\vb{x}}{\vb{A}\vb{x}} = \ip{\vb{x}}{\vb{A}^\top\vb{A}\vb{x}} = \ip{\vb{x}}{\vb{x}}
        \end{equation}
        which results in,
        \begin{equation}
          \begin{gathered}
            \vb{A}^\top\vb{A} = \vb{A}\vb{A}^\top = \vb{I} \\
            \vb{A}^{-1} = \vb{A}^\top
          \end{gathered}
        \end{equation}
        where $\vb{I}$ is the identity matrix. The inverse of an orthogonal transformation is
        the transpose of the matrix itself. As a note: if the matrix $\vb{A}$ had complex
        elements, then the transpose would be replaced by the conjugate transpose
        $\bar{A}^\top$:
        \begin{equation}
          \ip{\vb{A}\vb{x}}{\vb{A}\vb{x}} = \ip{\vb{x}}{\bar{A}^\top\vb{A}\vb{x}} = \ip{\vb{x}}{\vb{x}}
        \end{equation}

  \item \textbf{Orthogonality of rows and columns}: Considering the row and column vectors
        $\vb{v}_i$ and $\vb{a}_i$ forming the matrix $\vb{A}$, we can write the above
        equation as
        \begin{equation}
          \begin{gathered}
            \ip{\vb{v}_i}{\vb{v}_j} = \delta_{ij} \\
            \ip{\vb{a}_i}{\vb{a}_j} = \delta_{ij}
          \end{gathered}
        \end{equation}
        which is equivalent to requiring that the rows and the columns vectors are a set of
        orthonormal vectors. The vectors $\vb{v}_i$ represent the components of the unit
        vectors that define the direction of the coordinate axes in the rotated system,
        relative to the unit vectors of the reference system.

  \item \textbf{Property of eigenvalues}: The eigenvector $\vb{v}$ corresponding to the
        eigenvalue $\lambda$ will preserve its length under the orthogonal transformation:
        \begin{equation}
          \ip{\vb{v}}{\vb{v}} = \ip{\vb{A}\vb{v}}{\vb{A}\vb{v}} = \lambda\overline{\lambda}\ip{\vb{v}}{\vb{v}}
        \end{equation}
        which implies that the eigenvalues are complex numbers on the complex unit circle of
        the form $\lambda = e^{i\theta}$, and they always come in pairs, as
        $\lambda = -e^{i\theta}$ is an eigenvalue as well. If $n$ is even, then one eigenvalue
        is either $+1$ or $-1$.

  \item \textbf{Property of eigenvectors}: the eigenvectors $\vb{v}_1$ and $\vb{v}_2$
        corresponding to different eigenvalues $\lambda_1$ and $\lambda_2$ are orthogonal to
        each other. The proof of this property is straightforward, as we can write
        \begin{equation}
          \ip{\vb{v}_1}{\vb{v}_2} = \ip{\vb{A}\vb{v}_1}{\vb{A}\vb{v}_2} = \lambda_1\overline{\lambda_2}\ip{\vb{v}_1}{\vb{v}_2}
        \end{equation}
        which can only be true if $\ip{\vb{v}_1}{\vb{v}_2} = 0$, as the eigenvalues are
        complex numbers on the unit circle different from zero.

  \item \textbf{Property of determinant:} From $\det(\vb{A}) = \det(\vb{A}^\top)$ it follows
        that,
        \begin{equation}
          \det(\vb{A})^2 = 1
        \end{equation}
        which implies that the determinant of an orthogonal matrix is either $+1$ or $-1$.
        As we will see in the next section, the former corresponds to proper rotations, while
        the latter corresponds to improper rotations or reflections.

  \item \textbf{Interpretation of orthogonal transformation:} It is important to note that
        an orthogonal transformation can be seen in either of two ways:
        \begin{itemize}
          \item a \textbf{passive} transformation, where the components of the \textbf{same}
                vector are expressed in different bases;
          \item an \textbf{active} transformation, where a vector is actively transformed into
                a \textbf{new one} in the same reference system.
        \end{itemize}
        We will observe that in three dimensions, these two interpretations result in rotations
        that can be either clockwise or counterclockwise around the rotation axis.
\end{itemize}

% ============================================================
\section{Rotations in Three Dimensions}
% ============================================================

An orthogonal transformation in three dimensions will have the eigenvalues $+1$ or $-1$, and
$\lambda = e^{i\theta}$ and $\lambda = e^{-i\theta}$.

The eigenvectors corresponding to the eigenvalue $+1$ are the vectors that remain unchanged
under the transformation, forming the rotation axis. Conversely, the eigenvectors
corresponding to the eigenvalue $-1$ are the vectors that are reflected under the
transformation.

In physics, we assume that an orthogonal transformation varies continuously with the
parameter $\theta$ the rotation angle. This angle typically depends on time and describes the
evolution of the rotated reference system relative to the fixed axis. When the rotation angle
is zero, the reference systems coincide, and the transformation is the identity. Since
continuous transformations cannot reflect the rotation axis at any point in the dynamics, we
focus on proper rotations where the eigenvalue corresponding to the rotation axis is $+1$,
ensuring no reflections occur.

\subsection{Rotation Around the Principal Axes}

Before deriving the general form of a rotation around any axis, we will consider the rotation
around the $z$-axis. In \eqref{rot-figure} such a rotation is illustrated.
 The rotation matrix $\vb{A}_z(\theta)$ in this case is given by,

\begin{equation}
  \vb{A}_z(\theta) =
  \begin{pmatrix}
    \cos\theta  & \sin\theta & 0 \\
    -\sin\theta & \cos\theta & 0 \\
    0           & 0          & 1
  \end{pmatrix}
\end{equation}

where $\theta$ is the rotation angle. The derivation of the matrix $\vb{A}_z(\theta)$ is
straightforward, as the row vectors are obviously the components of the unit vectors of the
rotated system in the reference system. The matrix $\vb{A}_z(\theta)$ is orthogonal, as it
satisfies the properties of an orthogonal matrix, and it is easy to verify that the
determinant is $+1$. It is important to emphasize that the matrix $\vb{A}_z(\theta)$
represents a passive transformation, where the components of the same vector are expressed
in different bases, and the rotated reference system is turned by an angle $\theta$
\textbf{counterclockwise} around the $z$-axis.

When $\vb{A}_z(\theta)$ is interpreted in an \textbf{active} case, the rotation will be
\textbf{clockwise} around the $z$-axis. This can be readily verified by considering the
effect of applying $\vb{A}_z(\theta)$ to a vector
$\vb{v} = (r\cos\phi,\, r\sin\phi,\, z)$,
\begin{equation}
  \vb{A}_z(\theta)\,\vb{v} =
  \begin{pmatrix}
    \cos\theta  & \sin\theta & 0 \\
    -\sin\theta & \cos\theta & 0 \\
    0           & 0          & 1
  \end{pmatrix}
  \begin{pmatrix} r\cos\phi \\ r\sin\phi \\ z \end{pmatrix}
  =
  \begin{pmatrix} r\cos(\phi-\theta) \\ r\sin(\phi-\theta) \\ z \end{pmatrix}
\end{equation}
which corresponds to decreasing the polar angle $\phi$ by the angle $\theta$. If the active
rotation in practical applications needs to be counterclockwise, the angle $\theta$ must be
negative.

The rotation matrix around the $x$-axis is given by,
\begin{equation}
  \vb{A}_x(\theta) =
  \begin{pmatrix}
    1 & 0           & 0          \\
    0 & \cos\theta  & \sin\theta \\
    0 & -\sin\theta & \cos\theta
  \end{pmatrix}
\end{equation}
and the one around the $y$-axis is given by,
\begin{equation}
  \vb{A}_y(\theta) =
  \begin{pmatrix}
    \cos\theta  & 0 & \sin\theta \\
    0           & 1 & 0          \\
    -\sin\theta & 0 & \cos\theta
  \end{pmatrix}
\end{equation}

\begin{figure}[h]
  \centering
  \includegraphics[width=0.45\textwidth]{rotated.jpg}
  \caption{Fig~1: Rotation of an angle $\theta$ around the $z$ axis.}
  \label{rot-figure}
\end{figure}



\subsection{Rotation Around a General Axis}

The formula for a change of coordinates as a result of a rotation around an arbitrary axis
defined by the unit vector $\vb{n} = (n_x, n_y, n_z)$, with $\ip{\vb{n}}{\vb{n}} = 1$
and an angle $\theta$, is called the Rodrigues formula. While there are numerous
well-documented methods to derive this formula available in various resources, we will revisit
the derivation process, explaining the steps in detail.

Our goal is to find the components of the rotated unit vectors $\vb{e}'_i$ in the reference
system $\vb{e}_i$ as a function of the components of the unit vector $\vb{n}$ and the angle
$\theta$. These components will form the \textbf{rows} of the matrix describing the change of
coordinates.

A unit vector $\vb{e}_i$ in the fixed reference system can be decomposed into the component
$\vb{e}_i^{||}$ and $\vb{e}_i^{\perp}$ respectively parallel and perpendicular to $\vb{n}$:
\begin{align}
  \vb{e}_i &= \vb{e}_i^{||} + \vb{e}_i^{\perp} \\
  \vb{e}_i^{||} &= \ip{\vb{e}_i}{\vb{n}}\,\vb{n} \\
  \vb{e}_i^{\perp} &= \vb{e}_i - \vb{e}_i^{||} = \vb{e}_i - \ip{\vb{e}_i}{\vb{n}}\,\vb{n}
\end{align}

The component $\vb{e}_i^{||}$ is parallel to the rotation axis, and it will not change under
the rotation. The component $\vb{e}_i^{\perp}$ is perpendicular to the rotation axis, and it
will rotate by an angle $\theta$ around the axis $\vb{n}$, into a new vector
$\vb{e}_i^{'\perp}$. Considering the unit vector along
$\vb{n} \times \vb{e}_i^{\perp} = \vb{n} \times \vb{e}_i$,
\begin{equation}
  \frac{\vb{n} \times \vb{e}_i^{\perp}}{|\vb{n} \times \vb{e}_i^{\perp}|}
  = \frac{\vb{n} \times \vb{e}_i}{|\vb{n} \times \vb{e}_i|}
\end{equation}
and $\vb{e}_i^{\perp}$,
\begin{equation}
  \frac{\vb{e}_i^{\perp}}{|\vb{e}_i^{\perp}|}
  = \frac{\vb{e}_i - \ip{\vb{e}_i}{\vb{n}}\,\vb{n}}{\sqrt{1 - \ip{\vb{e}_i}{\vb{n}}^2}}
\end{equation}
utilizing the rotation formula around the axis $z$ (in this case $\vb{n}$) we can write the
rotated vector as,
\begin{equation}
  \vb{e}_i^{'\perp} =
  \left(
    \cos\theta\,\frac{\vb{e}_i - \ip{\vb{e}_i}{\vb{n}}}{\sqrt{1 - \ip{\vb{e}_i}{\vb{n}}^2}}
    - \sin\theta\,\frac{\vb{n} \times \vb{e}_i}{|\vb{n} \times \vb{e}_i|}
  \right)\sqrt{1 - \ip{\vb{e}_i}{\vb{n}}^2}
\end{equation}
as the norm of the vector $\vb{e}_i^{\perp}$ remains constant. The vector
$\vb{n} \times \vb{e}_i$ contains only the components of $\vb{n}$ of index different from
$i$, as it can be easily verified by the definition of the cross product,
\begin{equation}
  (\vb{n} \times \vb{e}_i)_j = \epsilon_{jkp}\, n_k\,(e_i)_p = \epsilon_{jkp}\, n_k\,\delta_{ip} = \epsilon_{jki}\, n_k
\end{equation}
where $\epsilon_{ijk}$ is the Levi-Civita tensor. The norm of the vector $\vb{n} \times \vb{e}_i$ is given by,
\begin{align}
  |\vb{n} \times \vb{e}_i| &= \sqrt{n_k^2 + n_p^2}, \quad k, p \neq i \\
                            &= \sqrt{1 - \ip{\vb{e}_i}{\vb{n}}^2}
\end{align}

The components of the rotated unit vectors $\vb{e}'_i$ with respect to the reference frame
are then given by,
\begin{equation}
  \vb{e}'_i = \ip{\vb{e}_i}{\vb{n}}\,\vb{n}
    + \cos\theta\bigl(\vb{e}_i - \ip{\vb{e}_i}{\vb{n}}\bigr)
    + \sin\theta\,(\vb{e}_i \times \vb{n})
\end{equation}

The matrix $\vb{A}(\vb{n}, \theta)$ which describes the transformation of a vector's
components with respect to a reference system rotated around an axis $\vb{n}$ by an angle
$\theta$, is obtained by evaluating the above equation for each unit vector $\vb{e}_i$ in
the reference system. The components of the vectors $\vb{e}'_i$ with respect to $\vb{e}_i$
are the rows of the matrix $\vb{A}(\vb{n}, \theta)$:
\begin{equation}
  \vb{A}(\vb{n},\theta) = \cos(\theta)\,\vb{I} + (1-\cos(\theta))\ket{\vb{n}}\bra{\vb{n}}
    + \sin(\theta)\,(\cdot) \times \vb{n}
  \label{eq:rotation_matrix}
\end{equation}
where $(\cdot)$ is the placeholder for a vector. The explicit matrix form of the equation
above is,
\begin{equation}
  \vb{A}(\vb{n},\theta) =
  \begin{pmatrix}
    n_x^2(1-\cos\theta)+\cos\theta        & n_x n_y(1-\cos\theta)+n_z\sin\theta   & n_x n_z(1-\cos\theta)-n_y\sin\theta \\
    n_y n_x(1-\cos\theta)-n_z\sin\theta   & n_y^2(1-\cos\theta)+\cos\theta        & n_y n_z(1-\cos\theta)+n_x\sin\theta \\
    n_z n_x(1-\cos\theta)+n_y\sin\theta   & n_z n_y(1-\cos\theta)-n_x\sin\theta   & n_z^2(1-\cos\theta)+\cos\theta
  \end{pmatrix}
\end{equation}

It is important to emphasize again that the above transformation represents a passive
rotation, meaning that the components of the same vector are expressed in different bases.
The same matrix, if interpreted as an active transformation, will rotate the vector around
the axis $\vb{n}$ by an angle $\theta$ in the counterclockwise direction.

% ============================================================
\section{Infinitesimal Rotations}
% ============================================================

\subsection{Definition}

In the previous section, we demonstrated that a rotation in three dimensions can be uniquely
identified by a unit vector $\vb{n}$ and an angle $\theta$. It would seem natural to consider
whether the group of orthogonal matrices, referred to as $SO(3)$, is isomorphic to
$\mathbb{R}^3$. Let us postulate the creation of such a mapping $f$, such that,
\begin{equation}
  f: SO(3) \to \mathbb{R}^3
\end{equation}
For this mapping to be an isomorphism, it must be one-to-one (and surjective),
\begin{equation}
  \forall\; \vb{R}_1, \vb{R}_2 \in SO(3),\quad
  f(\vb{R}_1) = f(\vb{R}_2) \implies \vb{R}_1 = \vb{R}_2.
\end{equation}

Considering two rotations $\vb{A} \in SO(3)$ and $\vb{B} \in SO(3)$, we know that rotations
do not commute,
\begin{equation}
  \vb{A}\vb{B} \neq \vb{B}\vb{A}
\end{equation}
However, their images under the mapping $f$ are the same,
\begin{equation}
  f(\vb{A}\vb{B}) = f(\vb{A}) + f(\vb{B}) = f(\vb{B}) + f(\vb{A}) = f(\vb{B}\vb{A})
\end{equation}
due to the commutativity of the vector space $\mathbb{R}^3$. Thus, it is evident that such a
mapping cannot be one-to-one, and the group $SO(3)$ is not isomorphic to $\mathbb{R}^3$.

Finite rotations cannot therefore be uniquely identified by a vector in $\mathbb{R}^3$, and
we must consider a different approach. We will now introduce the concept of infinitesimal
rotations, which are rotations that do commute.

Let us consider two rotation matrices $\vb{A}$ and $\vb{B}$ such that,
\begin{align}
  \vb{A} &= \vb{I} + \epsilon_1 \\
  \vb{B} &= \vb{I} + \epsilon_2
\end{align}
where $\epsilon_1$ and $\epsilon_2$ are infinitesimal matrices such that
$\epsilon_1\epsilon_2 = \epsilon_2\epsilon_1 \approx 0$. The product of $\vb{A}$ and $\vb{B}$
is,
\begin{equation}
  \vb{A}\vb{B} = (\vb{I}+\epsilon_1)(\vb{I}+\epsilon_2) = \vb{I} + \epsilon_1 + \epsilon_2 + \epsilon_1\epsilon_2 = \vb{I} + \epsilon_1 + \epsilon_2
\end{equation}
which is the same as the product of $\vb{B}$ and $\vb{A}$. Hence, for so-called infinitesimal
rotations, the commutativity property holds.

For a general infinitesimal rotation $\vb{A} = \vb{I} + \epsilon$, the unitary condition
$\vb{A}^\dagger\vb{A} = \vb{I}$ implies that the matrix $\epsilon$ is skew-symmetric, as it
satisfies the property
\begin{equation}
  \epsilon^\top = -\epsilon.
\end{equation}
A skew-symmetric matrix has zero diagonal elements, and the off-diagonal elements satisfy the
property $\epsilon_{ij} = -\epsilon_{ji}$, and assumes the general form,
\begin{equation}
  \epsilon =
  \begin{pmatrix}
    0           &  d\Omega_3 & -d\Omega_2 \\
    -d\Omega_3  &  0         &  d\Omega_1 \\
    d\Omega_2   & -d\Omega_1 &  0
  \end{pmatrix}
\end{equation}

A change of coordinates of a vector $\vb{r}$ under an infinitesimal rotation
$\vb{A} = \vb{I} + \epsilon$ is given by,
\begin{equation}
  \vb{r}' = \vb{A}\vb{r} = (\vb{I} + \epsilon)\vb{r} = \vb{r} + \epsilon\vb{r}
           = \vb{r} + \vb{r} \times \boldsymbol{d\Omega}
\end{equation}
as it can be readily verified that,
\begin{equation}
  \epsilon\vb{r} =
  \begin{pmatrix}
    0           &  d\Omega_3 & -d\Omega_2 \\
    -d\Omega_3  &  0         &  d\Omega_1 \\
    d\Omega_2   & -d\Omega_1 &  0
  \end{pmatrix}
  \begin{pmatrix} x \\ y \\ z \end{pmatrix}
  =
  \begin{pmatrix}
    d\Omega_3 y - d\Omega_2 z \\
    -d\Omega_3 x + d\Omega_1 z \\
    d\Omega_2 x - d\Omega_1 y
  \end{pmatrix}
  = \vb{r} \times d\boldsymbol{\Omega}
\end{equation}
where $\vb{r} = (x,y,z)$ and $\boldsymbol{\Omega} = (d\Omega_1, d\Omega_2, d\Omega_3)$.

As two infinitesimal rotations commute, the mapping $f$ from the group of infinitesimal
rotations to $\mathbb{R}^3$ is one-to-one and surjective, and we can therefore uniquely
identify a rotation by the vector $\boldsymbol{\Omega}$.

We now evaluate the form of the matrix \eqref{eq:rotation_matrix} for a value of the angle
$\theta$ that is infinitesimally small. We find that the matrix is given by,
\begin{equation}
  \vb{R}(\vb{n}, d\theta) = \vb{I} + d\theta
  \begin{pmatrix}
    0     &  n_3  & -n_y \\
    -n_z  &  0    &  n_x \\
    n_y   & -n_x  &  0
  \end{pmatrix}
\end{equation}
This allows us to identify,
\begin{equation}
  \epsilon = d\theta
  \begin{pmatrix}
    0     &  n_3  & -n_y \\
    -n_z  &  0    &  n_x \\
    n_y   & -n_x  &  0
  \end{pmatrix}
\end{equation}
as the infinitesimal rotation matrix, and the vector
\begin{equation}
  \boldsymbol{\Omega} = d\theta\,(n_x, n_y, n_z)
\end{equation}
as the vector defining the infinitesimal rotation.

\subsection{Similarity Transformation of Infinitesimal Rotations}

In the previous section, we have shown that an infinitesimal rotation can be represented by
a vector $\boldsymbol{\Omega} = (d\Omega_1, d\Omega_2, d\Omega_3)$. We will now show that
$\boldsymbol{\Omega}$ transforms as a regular $\mathbb{R}^3$ vector under a similarity
transformation.

Let us consider a rotation matrix $\vb{A} = \vb{I} + \epsilon$, where $\epsilon$ is a
skew-symmetric matrix of the form,
\begin{equation}
  \epsilon =
  \begin{pmatrix}
    0           &  d\Omega_3 & -d\Omega_2 \\
    -d\Omega_3  &  0         &  d\Omega_1 \\
    d\Omega_2   & -d\Omega_1 &  0
  \end{pmatrix}
\end{equation}

The representation of $\vb{A}$ in the coordinate system defined by an orthogonal
transformation $\mathbf{B}$, is $\vb{A}' = \mathbf{B}\vb{A}\mathbf{B}^\top$ and it has the
form,
\begin{equation}
  \vb{A}' = \mathbf{B}(\vb{I}+\epsilon)\mathbf{B}^\top
           = \mathbf{B}\mathbf{B}^\top + \mathbf{B}\epsilon\mathbf{B}^\top
           = \vb{I} + \mathbf{B}\epsilon\mathbf{B}^\top
           = I + \epsilon'
\end{equation}

The matrix $\epsilon' = \mathbf{B}\epsilon\mathbf{B}^\top$ is skew-symmetric, as it can be
readily verified that,
\begin{equation}
  (\epsilon')^\top = (\mathbf{B}\epsilon\mathbf{B}^\top)^\top
                   = \mathbf{B}\epsilon^\top\mathbf{B}^\top
                   = -\mathbf{B}\epsilon\mathbf{B}^\top = -\epsilon'
\end{equation}

Given that the matrix $\mathbf{B}$ is orthogonal, there exist three orthogonal unit vectors
$\mathbf{v}_1$, $\mathbf{v}_2$, and $\mathbf{v}_3$:
\begin{equation}
  \mathbf{v}_i \cdot \mathbf{v}_j = \delta_{ij}
\end{equation}
which define a reference system in the rotated coordinates. The matrix $\mathbf{B}$ can be
written as:
\begin{equation}
  \mathbf{B} =
  \begin{pmatrix} \mathbf{v}_1 \\ \mathbf{v}_2 \\ \mathbf{v}_3 \end{pmatrix}
\end{equation}
where the vectors $\mathbf{v}_i$ are row vectors. The transpose of $\mathbf{B}$ is readily
obtained as:
\begin{equation}
  \mathbf{B}^\top = \bigl[\mathbf{v}_1^\top,\; \mathbf{v}_2^\top,\; \mathbf{v}_3^\top\bigr]
\end{equation}
where the vectors $\mathbf{v}_i$ are now the columns of $\mathbf{B}^\top$.

The product $\boldsymbol{\epsilon}\mathbf{B}^\top$ is the matrix with columns:
\begin{equation}
  \boldsymbol{\epsilon}\mathbf{B}^\top =
  \bigl[\mathbf{v}_1 \times d\boldsymbol{\Omega},\;
        \mathbf{v}_2 \times d\boldsymbol{\Omega},\;
        \mathbf{v}_3 \times d\boldsymbol{\Omega}\bigr]
\end{equation}

and, for illustration, the matrix $\mathbf{B}\boldsymbol{\epsilon}\mathbf{B}^\top$ is
explicitly written below:
\begin{equation}
  \mathbf{B}\boldsymbol{\epsilon}\mathbf{B}^\top =
  \begin{pmatrix}
    \mathbf{v}_1\cdot(\mathbf{v}_1\times d\boldsymbol{\Omega}) &
    \mathbf{v}_1\cdot(\mathbf{v}_2\times d\boldsymbol{\Omega}) &
    \mathbf{v}_1\cdot(\mathbf{v}_3\times d\boldsymbol{\Omega}) \\
    \mathbf{v}_2\cdot(\mathbf{v}_1\times d\boldsymbol{\Omega}) &
    \mathbf{v}_2\cdot(\mathbf{v}_2\times d\boldsymbol{\Omega}) &
    \mathbf{v}_2\cdot(\mathbf{v}_3\times d\boldsymbol{\Omega}) \\
    \mathbf{v}_3\cdot(\mathbf{v}_1\times d\boldsymbol{\Omega}) &
    \mathbf{v}_3\cdot(\mathbf{v}_2\times d\boldsymbol{\Omega}) &
    \mathbf{v}_3\cdot(\mathbf{v}_3\times d\boldsymbol{\Omega})
  \end{pmatrix}
\end{equation}

The antisymmetry of the matrix is evident as the diagonal elements are zero:
\begin{equation}
  \mathbf{v}_i\cdot(\mathbf{v}_i\times d\boldsymbol{\Omega}) = 0
\end{equation}
and the off-diagonal elements are opposite because of the properties of the scalar triple
product:
\begin{equation}
  \mathbf{v}_i\cdot(\mathbf{v}_j\times d\boldsymbol{\Omega}) =
  -\mathbf{v}_j\cdot(\mathbf{v}_i\times d\boldsymbol{\Omega})
\end{equation}

The components of $d\boldsymbol{\Omega}'$ are by definition:
\begin{align}
  d\Omega'_1 &= \mathbf{v}_2\cdot(\mathbf{v}_3\times d\boldsymbol{\Omega})
              = d\boldsymbol{\Omega}\cdot(\mathbf{v}_2\times\mathbf{v}_3) \\
  d\Omega'_2 &= \mathbf{v}_3\cdot(\mathbf{v}_1\times d\boldsymbol{\Omega})
              = d\boldsymbol{\Omega}\cdot(\mathbf{v}_3\times\mathbf{v}_1) \\
  d\Omega'_3 &= \mathbf{v}_1\cdot(\mathbf{v}_2\times d\boldsymbol{\Omega})
              = d\boldsymbol{\Omega}\cdot(\mathbf{v}_1\times\mathbf{v}_2)
\end{align}

As the determinant of $\mathbf{B}$ assumes the values $\pm 1$:
\begin{equation}
  |\mathbf{B}| = \mathbf{v}_1\cdot(\mathbf{v}_2\times\mathbf{v}_3) = \pm 1
\end{equation}
depending on the orientation of the vectors $\mathbf{v}_i$ in the rotated coordinate system,
the equations become:
\begin{align}
  d\Omega'_1 &= |\mathbf{B}|\; d\boldsymbol{\Omega}\cdot\mathbf{v}_1 \\
  d\Omega'_2 &= |\mathbf{B}|\; d\boldsymbol{\Omega}\cdot\mathbf{v}_2 \\
  d\Omega'_3 &= |\mathbf{B}|\; d\boldsymbol{\Omega}\cdot\mathbf{v}_3
\end{align}
where we have made use of the property of the scalar triple product, or equivalently,
\begin{equation}
  d\boldsymbol{\Omega}' = |\mathbf{B}|\;\mathbf{B}\,d\boldsymbol{\Omega}.
\end{equation}

The vector $d\boldsymbol{\Omega}'$ transforms as a regular vector under a similarity
transformation, provided the orientation of the axis is preserved, and it is therefore a
vector in $\mathbb{R}^3$.

% ============================================================
\section{Rate of Change of a Vector}
% ============================================================

In this section, we will examine how the coordinates of a given vector change over time in
two different reference systems: one that is fixed and another that rotates with respect to
the fixed one. While there are numerous resources available that address this topic, we
believe it is essential to derive the fundamental equation for the rate of change in a
detailed manner. This will ensure that it is clear which components are being referenced in
the general equation of motion.

To understand this better, consider a vector in a fixed reference frame and how its
components might appear different in a rotating frame. The rotating frame itself can change
its orientation with time, adding complexity to the motion of the vector. By carefully
deriving the equations, we can delineate the contributions from both the intrinsic change of
the vector and the effect of the rotating frame.

We will start by defining the vector in the fixed frame and introduce the rotational matrix
that relates the rotating frame to the fixed frame. We will then derive the time derivative
of the vector in both frames and show how they are connected.

Let us consider a time-dependent vector $\vb{g} = \vb{g}(t)$ in a fixed reference frame,
and its components in a rotating frame $\vb{g}' = \vb{g}'(t)$. We will assume that the
components in the two different reference systems change as per a rotation matrix $\vb{A}(t)$,
which is also time-dependent:
\begin{equation}
  \vb{g}'(t) = \vb{A}(t)\,\vb{g}(t)
\end{equation}

The time derivative of the vector in the rotating frame is given by,
\begin{equation}
  \frac{d\vb{g}'(t)}{dt} = \frac{d\vb{A}(t)}{dt}\,\vb{g}(t) + \vb{A}(t)\,\frac{d\vb{g}(t)}{dt}
\end{equation}

and the time derivative of the matrix $\vb{A}(t)$ at time $t_0$ is formally defined as,
\begin{equation}
  \frac{d\vb{A}(t)}{dt} = \lim_{\Delta t \to 0}\frac{\vb{A}(t+\Delta t) - \vb{A}(t)}{\Delta t}.
\end{equation}

The operator $\vb{A}(t+\Delta t)$ is the composition of two rotations: one from the fixed
frame to the rotating frame at time $t$, and the other being the infinitesimal rotation from
the rotating frame at time $t + dt$:
\begin{equation}
  \vb{A}(t+\Delta t) = (\vb{I} + \epsilon)\,\vb{A}(t)
\end{equation}
and the derivative of the operator $\vb{A}(t)$ is then given by,
\begin{equation}
  \frac{d\vb{A}(t)}{dt}
  = \lim_{\Delta t\to 0}\frac{(\vb{I}+\epsilon)\vb{A}(t) - \vb{A}(t)}{\Delta t}
  = \lim_{\Delta t\to 0}\frac{\epsilon}{\Delta t}\,\vb{A}(t)
  = \vb{A}(t)\,(\cdot)\times\lim_{\Delta t\to 0}\frac{d\boldsymbol{\Omega}}{\Delta t}
  = \vb{A}(t)\,(\cdot)\times\boldsymbol{\omega}
\end{equation}
where,
\begin{equation}
  \boldsymbol{\omega} = \lim_{\Delta t\to 0}\frac{d\boldsymbol{\Omega}}{\Delta t}
\end{equation}
and $(\cdot)$ denotes the placeholder for a vector. $\boldsymbol{\omega}$ is the
\textbf{angular velocity}. The final formula for the rate of change of the vector in the
rotating frame is then given by the following equivalent equations:
\begin{align}
  \frac{d\vb{g}'(t)}{dt}
    &= \bigl(\vb{A}(t)\,\vb{g}(t)\bigr) \times \boldsymbol{\omega}
       + \vb{A}(t)\,\frac{d\vb{g}(t)}{dt} \nonumber \\
    &= \vb{g}'(t) \times \boldsymbol{\omega} + \vb{A}(t)\,\frac{d\vb{g}(t)}{dt} \nonumber \\
    &= \vb{A}(t)\!\left(\vb{g}(t) \times \bigl(\vb{A}(t)^\top\boldsymbol{\omega}\bigr)
       + \frac{d\vb{g}(t)}{dt}\right)
  \label{eq:rate_direct}
\end{align}
where in the last equation, the distributive property of matrix multiplication for orthogonal
matrices regarding the cross product has been utilized.

A similar formula can be derived for the rate of change of the vector in the fixed frame,
since $\vb{g}(t) = \vb{A}^\top(t)\,\vb{g}'(t)$. Multiplying \eqref{eq:rate_direct} by
$\vb{A}^\top(t)$ from the left, we obtain,
\begin{align}
  \frac{d\vb{g}(t)}{dt}
    &= \bigl(\vb{A}^\top(t)\boldsymbol{\omega}\bigr) \times \vb{g}(t)
       + \vb{A}^\top(t)\,\frac{d\vb{g}'(t)}{dt} \nonumber \\
    &= \vb{A}^\top(t)\!\left(\boldsymbol{\omega} \times \vb{g}'(t)
       + \frac{d\vb{g}'(t)}{dt}\right)
  \label{eq:rate_inverse}
\end{align}

\subsection{Poisson Formulas}

Let us consider a set of fixed unit vectors in the rotating frame $\vb{e}'_i$ and their
corresponding vectors $\vb{e}_i$ in the fixed frame. From equation \eqref{eq:rate_direct} we
have,
\begin{equation}
  0 = \vb{e}'_i \times \boldsymbol{\omega} + \vb{A}\,\frac{d\vb{e}_i}{dt}
\end{equation}
which gives immediately the Poisson formula,
\begin{equation}
  \frac{d\vb{e}'_i}{dt} = \boldsymbol{\omega} \times \vb{e}'_i
\end{equation}
where,
\begin{equation}
  \frac{d\vb{e}'_i}{dt} = \vb{A}\,\frac{d\vb{e}_i}{dt}
\end{equation}

\subsection{Components of the Angular Velocity}

Let us consider the Poisson formula for the unit vector $\vb{e}'_1$,
\begin{equation}
  \boldsymbol{\omega} \times \vb{e}'_1 = \frac{d\vb{e}'_1}{dt}
\end{equation}
Performing the scalar product on both sides by $\vb{e}'_2$ from the left, we obtain the
component of the angular velocity along the $\vb{e}'_3$ vector:
\begin{align}
  \vb{e}'_2 \cdot \bigl(\boldsymbol{\omega} \times \vb{e}'_1\bigr)
    &= \vb{e}'_2 \cdot \frac{d\vb{e}'_1}{dt} \\
  \boldsymbol{\omega} \cdot (\vb{e}'_1 \times \vb{e}'_2)
    &= \vb{e}'_2 \cdot \frac{d\vb{e}'_1}{dt} \\
  \boldsymbol{\omega} \cdot \vb{e}'_3
    &= \vb{e}'_2 \cdot \frac{d\vb{e}'_1}{dt}
\end{align}

Similar formulas can be derived for the other components of the angular velocity, and we
ultimately obtain:
\begin{align}
  \omega_1 &= \vb{e}'_2 \cdot \frac{d\vb{e}'_3}{dt} \\
  \omega_2 &= \vb{e}'_3 \cdot \frac{d\vb{e}'_1}{dt} \\
  \omega_3 &= \vb{e}'_1 \cdot \frac{d\vb{e}'_2}{dt}
\end{align}

and in vector form,
\begin{equation}
  \boldsymbol{\omega} =
    \left(\vb{e}'_2 \cdot \frac{d\vb{e}'_3}{dt}\right)\vb{e}'_1
  + \left(\vb{e}'_3 \cdot \frac{d\vb{e}'_1}{dt}\right)\vb{e}'_2
  + \left(\vb{e}'_1 \cdot \frac{d\vb{e}'_2}{dt}\right)\vb{e}'_3
  \label{eq:angular_velocity_components}
\end{equation}

% ============================================================
\section{Relevant Examples of Angular Velocity}
% ============================================================

In this section, we will present several examples of angular velocity in a rotating system.
We will start with the simple scenario of a rigid body rotating around a fixed axis. Following
that, we will examine the rotation of a system around an axis that changes its orientation
over time.

\subsection{Rotation around a Fixed Axis}

In this section, we will assume that a rigid body is rotating around a fixed axis $\vb{n}$
by an angle $\theta = \theta(t)$. The time derivative of the rotation matrix
$\vb{A}(\vb{n}, \theta(t))$ from \eqref{eq:rotation_matrix} is given by:
\begin{equation}
  \frac{d\vb{A}(\vb{n},\theta(t))}{dt} =
  \bigl(-\sin(\theta)\,\vb{I} + \sin(\theta)\ket{\vb{n}}\bra{\vb{n}} + \cos(\theta)\,(\cdot)\times\vb{n}\bigr)\,\dot\theta
\end{equation}

It can be readily verified that:
\begin{equation}
  \frac{d\vb{A}(\vb{n},\theta(t))}{dt} = \bigl(\vb{A}(\vb{n},\theta(t))\,(\cdot)\bigr) \times \boldsymbol{\omega}
\end{equation}
where the angular velocity in the rotating frame is given by:
\begin{equation}
  \boldsymbol{\omega} = \dot\theta\,\vb{n}.
  \label{eq:angular_velocity_fixed}
\end{equation}

It is straightforward to verify that the components of the angular velocity in the fixed
frame are the same as those in the rotating frame, as
$\vb{A}(\vb{n},\theta(t))\,\vb{n} = \vb{n}$.

\subsection{Rotation around a Rotating Axis}

The formula for the angular velocity when the axis of rotation is also changing in time can be
derived following the same steps as in the previous example. The calculations are intensive,
so we present the final result for the angular velocity in the rotating frame:
\begin{equation}
  \boldsymbol{\omega} = (1-\cos(\theta))\,\dot{\vb{n}}\times\vb{n}
    + \sin(\theta)\,\dot{\vb{n}} + \dot\theta\,\vb{n}
  \label{eq:angular_velocity_rotating}
\end{equation}

The formula above can be verified using symbolic computation software. It is recommended to
express $\vb{n}$ in polar coordinates to automatically satisfy the orthogonality condition
$\vb{n}\cdot\dot{\vb{n}} = 0$ and ensure the preservation of the modulus of the vector
$\vb{n}$.

% ============================================================
\section{Equations of Motion in a Rotating Frame}
% ============================================================

In this section, we will thoroughly derive the equations of motion for a body in a rotating
frame of reference with respect to a fixed one. We will limit our discussion to cases of
constant angular velocity. The equations of motion will be expressed within the rotating
frame, introducing the Coriolis and centrifugal forces that emerge in these scenarios.

We start from the equation giving the rate of change of a vector in the fixed frame
\eqref{eq:rate_inverse}. We compute the second time derivative of the position vector
$\vb{r}$ in the fixed frame, and we consider its components in the rotating frame $\vb{r}'$:
\begin{equation}
  \frac{d^2\vb{r}(t)}{dt^2} =
    \frac{d\vb{A}^\top}{dt}\!\left(\boldsymbol{\omega}\times\vb{r}' + \frac{d\vb{r}'}{dt}\right)
    + \vb{A}^\top\!\left(\boldsymbol{\omega}\times\frac{d\vb{r}'}{dt} + \frac{d^2\vb{r}'}{dt^2}\right)
\end{equation}

From the orthogonality of the rotation matrix, $\vb{A}^\top\vb{A} = \vb{I}$:
\begin{equation}
  \frac{d\vb{A}^\top}{dt}\,\vb{A} + \vb{A}\,\frac{d\vb{A}^\top}{dt} = 0
\end{equation}
we obtain:
\begin{equation}
  \frac{d\vb{A}^\top}{dt} = -\vb{A}^\top\frac{d\vb{A}}{dt}\,\vb{A}^\top
\end{equation}

Inserting this expression in the equation above, and using the distributive property of the
cross product for orthogonal matrices,
\begin{equation}
  \begin{aligned}
    \frac{d^2\vb{r}(t)}{dt^2}
    &= -\vb{A}^\top\frac{d\vb{A}}{dt}\,\vb{A}^\top
       \!\left(\boldsymbol{\omega}\times\vb{r}' + \frac{d\vb{r}'}{dt}\right)
       + \vb{A}^\top\!\left(\boldsymbol{\omega}\times\frac{d\vb{r}'}{dt} + \frac{d^2\vb{r}'}{dt^2}\right) \\
    &= -\vb{A}^\top\frac{d\vb{A}}{dt}
       \!\left(\vb{A}^\top\boldsymbol{\omega}\times\vb{A}^\top\vb{r}'
         + \vb{A}^\top\frac{d\vb{r}'}{dt}\right)
       + \vb{A}^\top\!\left(\boldsymbol{\omega}\times\frac{d\vb{r}'}{dt} + \frac{d^2\vb{r}'}{dt^2}\right) \\
    &= -\vb{A}^\top\!\left(\vb{A}\!\left(\vb{A}^\top\boldsymbol{\omega}\times\vb{A}^\top\vb{r}'
         + \vb{A}^\top\frac{d\vb{r}'}{dt}\right)\right)\times\boldsymbol{\omega}
       + \vb{A}^\top\!\left(\boldsymbol{\omega}\times\frac{d\vb{r}'}{dt} + \frac{d^2\vb{r}'}{dt^2}\right) \\
    &= -\vb{A}^\top\!\left(\boldsymbol{\omega}\times\vb{r}' + \frac{d\vb{r}'}{dt}\right)\times\boldsymbol{\omega}
       + \vb{A}^\top\!\left(\boldsymbol{\omega}\times\frac{d\vb{r}'}{dt} + \frac{d^2\vb{r}'}{dt^2}\right)
  \end{aligned}
\end{equation}

we arrive at the equation:
\begin{equation}
  \vb{A}\,\frac{d^2\vb{r}(t)}{dt^2} =
    \boldsymbol{\omega}\times(\boldsymbol{\omega}\times\vb{r}')
    + 2\boldsymbol{\omega}\times\frac{d\vb{r}'}{dt}
    + \frac{d^2\vb{r}'}{dt^2}.
\end{equation}

Defining the acceleration in the fixed frame as $\vb{a} = \frac{d^2\vb{r}(t)}{dt^2}$, and
the acceleration in the rotating frame as $\vb{a}' = \frac{d^2\vb{r}'(t)}{dt^2}$, we obtain
the equation of motion:
\begin{equation}
  \vb{a}' = -\boldsymbol{\omega}\times(\boldsymbol{\omega}\times\vb{r}')
    - 2\boldsymbol{\omega}\times\frac{d\vb{r}'}{dt}
    + \vb{A}\,\vb{a}
  \label{eq:eom_rotating}
\end{equation}

The centrifugal force $\vb{F}_\text{cf}$ on a body of mass $m$ is the term
\begin{equation}
  \vb{F}_\text{cf} = -m\boldsymbol{\omega}\times(\boldsymbol{\omega}\times\vb{r}')
    = -m\bigl(\boldsymbol{\omega}\ip{\boldsymbol{\omega}}{\vb{r}'} - \omega^2\vb{r}'\bigr)
  \label{eq:centrifugal}
\end{equation}
which has a magnitude proportional to the square of the angular velocity, and is directed
radially outward from the axis of rotation. The Coriolis force is the term dependent on the
body velocity,
\begin{equation}
  F_\text{cor} = -2m\boldsymbol{\omega}\times\frac{d\vb{r}'}{dt}
  \label{eq:coriolis}
\end{equation}
which arises from the movement of the body in the rotating frame and is perpendicular to the
velocity of the body.

\subsection{Lagrange Equations of Motion in a Rotating Frame}

In this section and the next, we will shorten the symbol of time derivative $da/dt$ to
$\dot{a}$ to simplify the equations. Further, we will assume that all quantities are in the
rotating frame, and $\vb{r}$ represents the position vector in such a frame.

The equations of motion in a rotating frame can be expressed using Lagrangian mechanics. The
Lagrangian is defined as the difference between the kinetic and potential energies of the
system:
\begin{equation}
  L = T - V
\end{equation}
where $T$ is the kinetic energy and $V$ is the potential of the forces in the system.

The centrifugal force admits a potential $V_\text{cf}$ calculated as:
\begin{equation}
  V_\text{cf} = \frac{m}{2}\bigl(-\ip{\vb{r}}{\vb{r}}\ip{\boldsymbol{\omega}}{\boldsymbol{\omega}}
    + \ip{\boldsymbol{\omega}}{\vb{r}}^2\bigr)
\end{equation}
thus,
\begin{equation}
  \vb{F}_\text{cf} = -\frac{\partial V_\text{cf}}{\partial\vb{r}}
\end{equation}

The Coriolis force, similarly to the Lorentz force in electromagnetism, admits a generalized
potential $V_\text{cor}$ determined as,
\begin{equation}
  V_\text{cor} = -m\ip{\boldsymbol{\omega}\times\vb{r}}{\dot{\vb{r}}}
\end{equation}
so that the Coriolis force derives from a generalized potential:
\begin{equation}
  \vb{F}_\text{cor} = \frac{d}{dt}\frac{\partial V_\text{cor}}{\partial\dot{\vb{r}}}
    - \frac{\partial V_\text{cor}}{\partial\vb{r}}
\end{equation}

Consequently, the Lagrangian of the system becomes:
\begin{equation}
  L = \frac{1}{2}m\ip{\dot{\vb{r}}}{\dot{\vb{r}}}
    + \frac{m}{2}\bigl(-\ip{\vb{r}}{\vb{r}}\ip{\boldsymbol{\omega}}{\boldsymbol{\omega}}
      + \ip{\boldsymbol{\omega}}{\vb{r}}^2\bigr)
    - m\ip{\boldsymbol{\omega}\times\vb{r}}{\vb{v}}
\end{equation}

The Lagrange equations of motion are given by:
\begin{equation}
  \frac{d}{dt}\frac{\partial L}{\partial\dot{\vb{r}}} - \frac{\partial L}{\partial\vb{r}} = 0
\end{equation}

The first term of the Lagrangian gives:
\begin{equation}
  \frac{d}{dt}\frac{\partial L}{\partial\dot{\vb{r}}} = m\dot{\vb{r}}
\end{equation}

The second term of the Lagrangian gives:
\begin{equation}
  \frac{\partial L}{\partial\vb{r}} = -m\boldsymbol{\omega}\times(\boldsymbol{\omega}\times\vb{r})
    - 2m\boldsymbol{\omega}\times\dot{\vb{r}}
\end{equation}

The Lagrange equations of motion in a rotating frame are then:
\begin{equation}
  \ddot{\vb{r}} = -\boldsymbol{\omega}\times(\boldsymbol{\omega}\times\vb{r})
    - 2\boldsymbol{\omega}\times\dot{\vb{r}}
\end{equation}

It has to be noted that a conservative force in the inertial frame is not necessarily
conservative in the rotating frame. A simple example can be easily constructed by considering
a constant force in the plane $z=0$ of the inertial frame such as:
\begin{equation}
  \vb{F} = F_x\,\vb{i} + F_y\,\vb{j}
\end{equation}
where $F_x$ and $F_y$ are constants and $\vb{i}$, $\vb{j}$ are the unit vectors of the
inertial frame. The force in a frame rotating with constant angular velocity $\omega$ around
the axis $z$ is given by:
\begin{equation}
  \vb{F}' =
  \begin{pmatrix}
    \cos\omega t  & \sin\omega t & 0 \\
    -\sin\omega t & \cos\omega t & 0 \\
    0             & 0            & 1
  \end{pmatrix}
  \vb{F}
  = (F_x\cos(\omega t) + F_y\sin(\omega t))\,\vb{i}'
  + (-F_x\sin(\omega t) + F_y\cos(\omega t))\,\vb{j}'
\end{equation}
where $\vb{i}'$, $\vb{j}'$ are the unit vectors of the rotating frame. The force in the
rotating frame is time-dependent, and does not have a generalized potential as it does not
explicitly depend on the velocity. Each case will have to be considered individually to
determine whether the dynamics can be described in a Lagrangian formalism, and in general the
equations of motion \eqref{eq:eom_rotating} will have to be used.

% ============================================================
\section{Centrifugal and Coriolis Forces on Earth}
% ============================================================

In this section, we will delve into the Coriolis and centrifugal forces near Earth's surface.
However, we will significantly simplify the discussion by assuming that Earth is a perfect
sphere. Our goal here is to explore the mathematical formalism of the Coriolis and centrifugal
forces in an idealized scenario, providing a clear starting point for those interested in the
more realistic case.

We will focus on a reference frame fixed to the Earth, where the $z$ axis runs from the South
Pole to the North Pole, with positive direction from South to North. Earth's angular velocity
is along the $z$ axis:
\begin{equation}
  \boldsymbol{\omega} = \omega\,\vb{k}
\end{equation}
where $\omega$ is the angular velocity of the Earth equal to $7.29\times10^{-5}\,\text{s}^{-1}$.

We will be interested in the motion of a body in the vicinity of the Earth surface, and the
ideal coordinates to describe such a motion are the spherical coordinates $(\rho, \theta,
\phi)$, such that,
\begin{align}
  x &= \rho\sin\phi\cos\theta \\
  y &= \rho\sin\phi\sin\theta \\
  z &= \rho\cos\phi
\end{align}
where $r \geq 0$ is the radial distance from the center of the reference system,
$0 \leq \theta \leq 2\pi$ is the azimuthal angle, and $0 \leq \phi \leq \pi$ is the polar
angle. Associated with the spherical coordinates are the unit vectors $\vb{e}_r$,
$\vb{e}_\theta$, and $\vb{e}_\phi$ given by,
\begin{align}
  \vb{e}_\rho   &= \sin\phi\cos\theta\,\vb{i} + \sin\phi\sin\theta\,\vb{j} + \cos\phi\,\vb{k} \\
  \vb{e}_\theta &= -\sin\theta\,\vb{i} + \cos\theta\,\vb{j} \\
  \vb{e}_\phi   &= \cos\phi\cos\theta\,\vb{i} + \cos\phi\sin\theta\,\vb{j} - \sin\phi\,\vb{k}
\end{align}
where $\vb{i}$, $\vb{j}$, and $\vb{k}$ are the unit vectors of the rotating reference system
along the $x$, $y$, and $z$ axes, respectively.

The unit vectors define a local orthonormal basis and maintain the orientation of the
reference system:
\begin{align}
  \vb{e}_\theta \times \vb{e}_\rho &= \vb{e}_\phi \\
  \vb{e}_\phi   \times \vb{e}_\theta &= \vb{e}_\rho \\
  \vb{e}_\rho   \times \vb{e}_\phi &= \vb{e}_\theta
\end{align}
and they are related to the latitude and longitude of the Earth:
\begin{itemize}
  \item $\vb{e}_\rho$: Radial unit vector pointing outward from the center of the Earth,
        through the observer's position.
  \item $\vb{e}_\theta$: Points tangentially along lines of latitude towards the east.
  \item $\vb{e}_\phi$: Points tangentially along lines of longitude towards the South Pole.
\end{itemize}

The position vector in spherical coordinates is,
\begin{equation}
  \vb{r} = \rho\,\vb{e}_\rho
\end{equation}
and the angular velocity is represented by,
\begin{equation}
  \boldsymbol{\omega} = \omega\bigl(\cos(\phi)\,\vb{e}_\rho - \sin(\phi)\,\vb{e}_\phi\bigr).
\end{equation}
and finally the velocity of the body in spherical coordinates is given by,
\begin{equation}
  \dot{\vb{r}} = \dot\rho\,\vb{e}_\rho + \rho\dot\theta\sin(\phi)\,\vb{e}_\theta
    + \rho\dot\phi\,\vb{e}_\phi
\end{equation}

\subsection{Forces in Spherical Coordinates}

In the rotating frame there are three forces acting on the body: the gravitational force, the
centrifugal force, and the Coriolis force.

\begin{itemize}
  \item The gravitational force is directed radially inward and is given by,
  \begin{equation}
    \vb{F}_\text{g} = -mg\,\vb{e}_\rho
  \end{equation}
  where $m$ is the mass of the body and $g$ is the acceleration due to gravity.

  \item The centrifugal force is directed radially outward from the axis of rotation and in
        polar coordinates it is represented by the formula:
  \begin{equation}
    \vb{F}_\text{cf} = m\omega^2\rho\bigl(\sin^2(\phi)\,\vb{e}_\rho
      + \sin(\phi)\cos(\phi)\,\vb{e}_\phi\bigr)
    \label{eq:centrifugal_spherical}
  \end{equation}

  \item The Coriolis force is perpendicular to the velocity of the body and is quantified by
        the formula:
  \begin{align}
    \vb{F}_\text{cor}
    &= -2m\boldsymbol{\omega}\times\dot{\vb{r}}
     = -2m\omega\bigl(\cos(\phi)\,\vb{e}_\rho - \sin(\phi)\,\vb{e}_\phi\bigr)
       \times\bigl(\dot\rho\,\vb{e}_\rho + \rho\dot\theta\sin(\phi)\,\vb{e}_\theta + \rho\dot\phi\,\vb{e}_\phi\bigr) \nonumber \\
    &= -2m\omega\bigl(\rho\dot\theta\sin(\phi)\cos(\phi)\,\vb{e}_\rho\times\vb{e}_\theta
       + (\cos(\phi)\rho\dot\phi + \sin(\phi)\dot\rho)\,\vb{e}_\rho\times\vb{e}_\phi
       - \sin^2(\phi)\rho\dot\theta\,\vb{e}_\phi\times\vb{e}_\theta\bigr) \nonumber \\
    &= -2m\omega\bigl(-\rho\dot\theta\sin(\phi)\cos(\phi)\,\vb{e}_\phi
       + (\cos(\phi)\rho\dot\phi + \sin(\phi)\dot\rho)\,\vb{e}_\theta
       - \sin^2(\phi)\rho\dot\theta\,\vb{e}_\rho\bigr)
  \end{align}
  which gives the final formula,
  \begin{equation}
    \vb{F}_\text{cor} = m\omega\bigl(
      2\rho\dot\theta\sin^2(\phi)\,\vb{e}_\rho
      - 2(\dot\rho\sin(\phi) + \rho\dot\phi\cos(\phi))\,\vb{e}_\theta
      + \rho\dot\theta\sin(2\phi)\,\vb{e}_\phi\bigr)
    \label{eq:coriolis_spherical}
  \end{equation}
\end{itemize}

Before proceeding, it is interesting to consider the total force acting on the body in the
rotating frame without the gravitational force,
\begin{equation}
  \begin{aligned}
    \vb{F} &= m\omega\rho(t)\sin^2(\phi(t))\bigl(2\dot\theta(t)+\omega\bigr)\,\vb{e}_\rho \\
    &\quad - 2m\omega\bigl(\dot\rho(t)\sin(\phi(t)) + \rho(t)\dot\phi(t)\cos(\phi(t))\bigr)\,\vb{e}_\theta \\
    &\quad + m\omega\rho(t)\sin(\phi(t))\cos(\phi(t))\bigl(2\dot\theta(t)+\omega\bigr)\,\vb{e}_\phi
  \end{aligned}
\end{equation}

We now consider two cases, one where the body is at rest in the inertial frame, and one where
the total force is zero in the rotating frame.

In the first case, the body in the rotating frame will appear to have an angular velocity
$\dot\theta = -\omega$, and the other two velocities $\dot\rho$ and $\dot\phi$ are zero.
Under these conditions, the total force acting on the body is,
\begin{equation}
  \vb{F} = -m\omega^2\rho\bigl(\sin^2(\phi)\,\vb{e}_\rho + \sin(\phi)\cos(\phi)\,\vb{e}_\phi\bigr)
\end{equation}
which is a centripetal force equal to exactly the centrifugal force: the Coriolis force is the
one that counterbalances the centrifugal force and is responsible for the rotation in the
non-inertial frame.

In the second case, it is readily found that it is sufficient that $\dot\rho = 0$ and
$\dot\phi = 0$ and $\dot\theta = -\omega/2$, for the total force to be zero. Under these
conditions, the body in the rotating frame will appear to rotate counterclockwise with an
angular velocity equal to half the angular velocity of the rotating frame.

\subsection{Lagrangian in Spherical Coordinates}

We will now dive into the formulation of the Lagrangian of the system in spherical
coordinates. The potential of the gravitational force in spherical coordinates is given by,
\begin{equation}
  V_g = mg\,\rho(t)
\end{equation}
the potential of the centrifugal force is,
\begin{equation}
  V_\text{cf} = -\frac{1}{2}m\omega^2\rho(t)^2\sin^2(\phi(t))
\end{equation}
and finally the potential of the Coriolis force is represented by,
\begin{equation}
  V_\text{cor} = -m\omega\rho(t)^2\dot\theta(t)\sin^2(\phi(t)).
\end{equation}

The potential of the Coriolis force is the generalized potential when expressed in spherical
coordinates, such that,
\begin{align}
  \frac{d}{dt}\frac{\partial V_\text{cor}}{\partial\dot\theta}
    - \frac{\partial V_\text{cor}}{\partial\theta}
    &= \left\langle\vb{F}_\text{cor}\,\middle|\,\frac{\partial\vb{r}}{\partial\theta}\right\rangle \\
  \frac{d}{dt}\frac{\partial V_\text{cor}}{\partial\dot\phi}
    - \frac{\partial V_\text{cor}}{\partial\phi}
    &= \left\langle\vb{F}_\text{cor}\,\middle|\,\frac{\partial\vb{r}}{\partial\phi}\right\rangle \\
  \frac{d}{dt}\frac{\partial V_\text{cor}}{\partial\dot\rho}
    - \frac{\partial V_\text{cor}}{\partial\rho}
    &= \left\langle\vb{F}_\text{cor}\,\middle|\,\frac{\partial\vb{r}}{\partial\rho}\right\rangle
\end{align}

and similar equations hold for the centrifugal force,
\begin{align}
  -\frac{\partial V_\text{cf}}{\partial\theta}
    &= \left\langle\vb{F}_\text{cf}\,\middle|\,\frac{\partial\vb{r}}{\partial\theta}\right\rangle \\
  -\frac{\partial V_\text{cf}}{\partial\phi}
    &= \left\langle\vb{F}_\text{cf}\,\middle|\,\frac{\partial\vb{r}}{\partial\phi}\right\rangle \\
  -\frac{\partial V_\text{cf}}{\partial\rho}
    &= \left\langle\vb{F}_\text{cf}\,\middle|\,\frac{\partial\vb{r}}{\partial\rho}\right\rangle
\end{align}

The kinetic energy reads,
\begin{equation}
  T = \frac{1}{2}m\left(\rho(t)^2\bigl(\dot\theta(t)^2\sin^2(\phi(t)) + \dot\phi(t)^2\bigr)
    + \dot\rho(t)^2\right)
\end{equation}
and the Lagrangian of the system is,
\begin{equation}
  L = \frac{1}{2}m\left(-2g\rho(t) + \rho(t)^2\left(\sin^2(\phi(t))\bigl(\dot\theta(t)+\omega\bigr)^2
    + \dot\phi(t)^2\right) + \dot\rho(t)^2\right).
\end{equation}

The full equations of motion can be derived from the Euler--Lagrange equations,
\begin{align}
  \frac{d}{dt}\frac{\partial L}{\partial\dot\rho} - \frac{\partial L}{\partial\rho} &= 0 \\
  \frac{d}{dt}\frac{\partial L}{\partial\dot\phi} - \frac{\partial L}{\partial\phi} &= 0 \\
  \frac{d}{dt}\frac{\partial L}{\partial\dot\theta} - \frac{\partial L}{\partial\theta} &= 0
\end{align}
which result in the system of non-linear ordinary differential equations,
\begin{equation}
  \begin{gathered}
    g + \rho(t)\left(-\sin^2(\phi(t))\bigl(\dot\theta(t)+\omega\bigr)^2 - \dot\phi(t)^2\right)
      + \ddot\rho(t) = 0 \\
    \rho(t)\left(2\ddot\phi(t) - \sin(2\phi(t))\bigl(\dot\theta(t)+\omega\bigr)^2\right)
      + 4\dot\rho(t)\dot\phi(t) = 0 \\
    \rho(t)\ddot\theta(t)\sin(\phi(t))
      + 2\bigl(\dot\theta(t)+\omega\bigr)\bigl(\dot\rho(t)\sin(\phi(t))
        + \rho(t)\dot\phi(t)\cos(\phi(t))\bigr) = 0
  \end{gathered}
  \label{eq:full_eom_polar}
\end{equation}

\subsection{Linearized Equations of Motion}

We consider now the motion of an object covering distances over the Earth surface much
smaller than Earth radius $R$ in any direction. We will assume the dynamics is in a
neighbourhood of the position identified by spherical coordinates $(\rho_0, \theta_0, \phi_0)$,
\begin{equation}
  \vb{r}_0 = \rho_0\,\vb{e}_{\rho_0}
\end{equation}
where $\rho_0 \approx R$.

Around this position a local system of reference can be constructed by the unit vectors
$\vb{e}_{\rho_0}$, $\vb{e}_{\theta_0}$, and $\vb{e}_{\phi_0}$ as,
\begin{align}
  \vb{e}_{\rho_0}   &= \sin\phi_0\cos\theta_0\,\vb{i} + \sin\phi_0\sin\theta_0\,\vb{j} + \cos\phi_0\,\vb{k} \\
  \vb{e}_{\theta_0} &= -\sin\theta_0\,\vb{i} + \cos\theta_0\,\vb{j} \\
  \vb{e}_{\phi_0}   &= \cos\phi_0\cos\theta_0\,\vb{i} + \cos\phi_0\sin\theta_0\,\vb{j} - \sin\phi_0\,\vb{k}
\end{align}

and a point in the vicinity of the position $\vb{r}_0$ has spherical coordinates
$(\rho_0+\delta\rho,\;\theta_0+\delta\theta,\;\phi_0+\delta\phi)$ and position,
\begin{equation}
  \vb{r} = (\rho_0+\delta\rho)\,\vb{e}_{\rho_0+\delta\rho}
\end{equation}
which up to first order reads,
\begin{equation}
  \vb{r} = \vb{r}_0
    + \delta\rho\,\vb{e}_{\rho_0}
    + \rho_0\sin(\phi_0)\,\delta\theta\,\vb{e}_{\theta_0}
    + \rho_0\,\delta\phi\,\vb{e}_{\phi_0}.
\end{equation}

Under the assumption that $\|\vb{r} - \vb{r}_0\| \ll R$, we have for the displacements in
each direction,
\begin{align}
  \delta\rho &\ll \rho_0 \\
  \rho_0\,\delta\phi &\ll \rho_0 \quad\text{or}\quad \delta\phi \ll 1 \\
  \rho_0\sin(\phi_0)\,\delta\theta &\ll \rho_0 \quad\text{or}\quad \sin(\phi_0)\,\delta\theta \ll 1
\end{align}

We will therefore assume that the dynamics can be described by spherical coordinates
\begin{equation}
  \bigl(\rho_0 + \rho(t),\quad \theta_0 + \theta(t),\quad \phi_0 + \phi(t)\bigr)
\end{equation}
where at all times $\rho(t) \ll \rho_0$, $\phi(t) \ll 1$, and
$\sin(\phi_0)\,\theta(t) \ll 1$.

The equations of motion \eqref{eq:full_eom_polar} can be expressed in terms of the
displacements $(\rho(t), \theta(t), \phi(t))$ and simplified to first order as,
\begin{equation}
  \begin{alignedat}{6}
    \ddot\rho(t) &=& -g
      &+\omega^2\sin^2(\phi_0)\,\rho(t)
      &+\omega^2\rho_0\sin(2\phi_0)\,\phi(t)
      &+2\omega\rho_0\sin^2(\phi_0)\,\dot\theta(t) \\
    2\rho_0\ddot\phi(t) &=& \omega^2\rho_0\sin(2\phi_0)
      &+\omega^2\sin(2\phi_0)\,\rho(t)
      &+2\omega^2\rho_0\cos(2\phi_0)\,\phi(t)
      &+2\omega\rho_0\sin(2\phi_0)\,\dot\theta(t) \\
    \rho_0\sin(\phi_0)\ddot\theta(t) &=&
      &-2\omega\sin(\phi_0)\,\dot\rho(t)
      &-2\omega\rho_0\cos(\phi_0)\,\dot\phi(t)
  \end{alignedat}
  \label{eq:linearized_eom}
\end{equation}

The Earth's angular velocity and radius have values,
\begin{align}
  \omega &= 7.29\times10^{-5}\,\text{s}^{-1} \\
  R      &= 6371\,\text{km}
\end{align}
and the term $\omega^2\rho_0\sin^2(\phi_0)$ assumes the value,
\begin{equation}
  \omega^2\rho_0\sin^2(\phi_0) = 0.0339\,\sin^2(\phi_0)\;\text{m/s}^2
\end{equation}
which is orders of magnitude smaller than the gravitational acceleration $g = 9.81\,\text{m/s}^2$
and can be neglected, simplifying to:
\begin{equation}
  \begin{alignedat}{6}
    \ddot\rho(t) &=& -g
      &+\omega^2\sin^2(\phi_0)\,\rho(t)
      &+\omega^2\rho_0\sin(2\phi_0)\,\phi(t)
      &+2\omega\rho_0\sin^2(\phi_0)\,\dot\theta(t) \\
    2\rho_0\ddot\phi(t) &=& \omega^2\rho_0\sin(2\phi_0)
      &+\omega^2\sin(2\phi_0)\,\rho(t)
      &+2\omega^2\rho_0\cos(2\phi_0)\,\phi(t)
      &+2\omega\rho_0\sin(2\phi_0)\,\dot\theta(t) \\
    \rho_0\sin(\phi_0)\ddot\theta(t) &=&
      &-2\omega\sin(\phi_0)\,\dot\rho(t)
      &-2\omega\rho_0\cos(\phi_0)\,\dot\phi(t)
  \end{alignedat}
\end{equation}

Defining dimensional variables of length as,
\begin{align}
  x &= \rho_0\,\phi & y &= \rho_0\sin(\phi_0)\,\theta & z &= \rho
\end{align}
and defining a local coordinate system as,
\begin{align}
  \vb{e}_z &= \vb{e}_{\rho_0} &
  \vb{e}_y &= \vb{e}_{\theta_0} &
  \vb{e}_x &= \vb{e}_{\phi_0}
\end{align}
the linearized equations of motion can be expressed as,
\begin{equation}
  \begin{alignedat}{6}
    \ddot{z} &=& -g
      &+\omega^2\sin^2(\phi_0)\,z
      &+\omega^2\sin(2\phi_0)\,x
      &+2\omega\sin(\phi_0)\,\dot{y} \\
    2\ddot{x} &=& \omega^2\rho_0\sin(2\phi_0)
      &+\omega^2\sin(2\phi_0)\,z
      &+2\omega^2\cos(2\phi_0)\,x
      &+4\omega\cos(\phi_0)\,\dot{y} \\
    \ddot{y} &=&
      &-2\omega\sin(\phi_0)\,\dot{z}
      &-2\omega\cos(\phi_0)\,\dot{x}
  \end{alignedat}
  \label{eq:linearized_eom_xyz}
\end{equation}

\subsection{Foucault Pendulum}

In this section, we will derive the equations of motion for the Foucault pendulum, applying
several simplifying assumptions to streamline the process. The Foucault pendulum is a
straightforward experiment that illustrates the Earth's rotation. It consists of a heavy mass
$m$ suspended at a latitude $\phi_0$ by a long wire or rod of length $l$, free to oscillate
in any direction. Over time, the plane of the pendulum's oscillation rotates due to the
Coriolis force acting upon it.

To derive the equations of motion for the Foucault pendulum, we begin with the equations
\eqref{eq:linearized_eom_xyz}. We then disregard the $\omega^2$ terms caused by the
centrifugal force, as they are significantly smaller than the terms involving $\omega$
generated by the Coriolis force. Under this assumption, we arrive at the equations:
\begin{equation}
  \begin{gathered}
    \ddot{z} = -g + 2\omega\sin(\phi_0)\,\dot{y} \\
    \ddot{x} = \omega\cos(\phi_0)\,\dot{y} \\
    \ddot{y} = -2\omega\sin(\phi_0)\,\dot{z} - 2\omega\cos(\phi_0)\,\dot{x}
  \end{gathered}
\end{equation}

From the equations above we can derive the total force acting on the body as,
\begin{equation}
  \vb{F}_\text{total} = m
  \begin{pmatrix} \ddot{z} \\ \ddot{x} \\ \ddot{y} \end{pmatrix}
  = m
  \begin{pmatrix}
    -g + 2\omega\sin(\phi_0)\,\dot{y} \\
    \omega\cos(\phi_0)\,\dot{y} \\
    -2\omega\sin(\phi_0)\,\dot{z} - 2\omega\cos(\phi_0)\,\dot{x}
  \end{pmatrix}
\end{equation}
the Coriolis component of the total force is,
\begin{equation}
  \vb{F}_\text{cor} = m
  \begin{pmatrix}
    2\omega\sin(\phi_0)\,\dot{y} \\
    \omega\cos(\phi_0)\,\dot{y} \\
    -2\omega\sin(\phi_0)\,\dot{z} - 2\omega\cos(\phi_0)\,\dot{x}
  \end{pmatrix}
\end{equation}
which has a generalized potential,
\begin{equation}
  V_\text{Cor} = -m\ip{\boldsymbol{\omega}\times\vb{p}}{\dot{\vb{p}}}.
\end{equation}

In the equation above the angular velocity $\boldsymbol{\omega}$ of the Earth is considered
in this system of reference and is given by,
\begin{equation}
  \boldsymbol{\omega} = \omega(-\sin(\phi_0)\,\vb{e}_x + \cos(\phi_0)\,\vb{e}_z)
\end{equation}
and the position vector $\vb{p}$ is given by $\vb{p} = x\,\vb{e}_x + y\,\vb{e}_y + z\,\vb{e}_z$.

The Lagrangian of the system, including the gravitational force reads,
\begin{equation}
  \begin{aligned}
    L &= \frac{1}{2}m\bigl(\dot{x}(t)^2 + \dot{y}(t)^2 + \dot{z}(t)^2\bigr) \\
      &\quad - m\omega\Bigl(y(t)\bigl(\cos(\phi_0)\,\dot{x}(t) + \sin(\phi_0)\,\dot{z}(t)\bigr)
        - \dot{y}(t)\bigl(\cos(\phi_0)\,x(t) + \sin(\phi_0)\,z(t)\bigr)\Bigr) + mgz(t)
  \end{aligned}
\end{equation}

Utilizing a local system of spherical coordinates with constant radius $(\rho = l, \theta, \phi)$
in the system $(x,y,z)$, we can express the Lagrangian as,
\begin{equation}
  \begin{aligned}
    L = \frac{1}{2}lm\Bigl(l\Bigl(&\dot\theta(t)\bigl(\sin^2(\phi(t))\bigl(2\omega\cos(\phi_0)
        + \dot\theta(t) + \omega\sin(\phi_0)\cos(\theta(t))\sin(2\phi(t))\bigr) \\
      &+2\omega\sin(\phi_0)\sin(\theta(t))\dot\phi(t) + \dot\phi(t)^2\bigr) - 2g\cos(\phi(t))\Bigr)\Bigr).
  \end{aligned}
\end{equation}

The Euler equations of motion derived from the Lagrangian above read,
\begin{align}
  \phi(t)\dot\phi(t)\bigl(\omega\cos(\phi_0) + \dot\theta(t)\bigr) &= 0 \\
  l\phi(t)\dot\theta(t)\bigl(2\omega\cos(\phi_0) + \dot\theta(t)\bigr)
    &= l\ddot\phi(t) + g\phi(t)
\end{align}
which are equivalent to,
\begin{align}
  \ddot\phi(t) &= \left(-\frac{g}{l} - \omega^2\cos^2(\phi_0)\right)\phi(t) \\
  \dot\theta(t) &= -\omega\cos(\phi_0).
\end{align}

The equation in $\phi(t)$ is a simple harmonic oscillator with a frequency of
$\omega_\phi = \sqrt{g/l}$, if the higher order terms in $\omega^2\cos^2(\phi_0)$ are
removed. The equation of motion for $\theta(t)$ is a constant angular velocity which
represents the precession of the Foucault pendulum. In the northern hemisphere, the
precession is counterclockwise as $\phi < \pi/2$, while in the southern hemisphere it is
clockwise. The period of the precession is given by,
\begin{equation}
  \tau = \frac{2\pi}{\omega|\cos(\phi_0)|}
    = \frac{24}{|\cos(\phi_0)|}\,\text{h}
    = \frac{24}{\sin(\phi'_0)}\,\text{h}
\end{equation}
where the last equation is obtained by substituting $\omega = 2\pi/24\,\text{h}$, and
$\phi'_0$ is the geographical latitude of the pendulum. The maximum value of the period is
at the North and South poles; it decreases when the pendulum approaches the equator where the
period is infinite.

\end{document}
